% ============================================================
%  6.2 Fundamentos de Ciencia de Datos   (capítulo 6, Marco teórico)
%  Se incluye desde secciones/07_marco_teorico.tex con \input.
% ============================================================
\section{Fundamentos de Ciencia de Datos}
\label{sec:fundamentos-cd}

\subsection{Ciencia de Datos y ciencia de datos geoespacial}
\label{subsec:ciencia-datos}

La Ciencia de Datos puede definirse como el campo que se ocupa de extraer conocimiento útil
de los datos mediante la integración de la estadística, la computación y el conocimiento del
dominio. \textcite{donoho2017}, en su revisión de cincuenta años de la disciplina, la
caracteriza como una ampliación de la estadística que abarca la recolección, preparación y
exploración de datos, su representación y transformación, la computación, el modelado, la
visualización y la reflexión sobre la propia práctica científica con datos.
\textcite{provost2013} subrayan, por su parte, que la Ciencia de Datos se orienta a la toma
de decisiones basada en datos, y que su valor no reside en los algoritmos en sí, sino en su
capacidad de responder preguntas bien formuladas con evidencia empírica.

Cuando los datos tienen una referencia espacial explícita, el análisis debe incorporar los
principios de la ciencia de la información geográfica (GIScience), término que
\textcite{goodchild1992} propuso para designar el estudio de las cuestiones fundamentales que
plantea el uso de la información geográfica, como la representación, la escala, la
incertidumbre y la dependencia espacial. \textcite{singleton2021} han denominado ciencia de
datos geográfica a la convergencia entre ambas tradiciones, que combina las herramientas
computacionales de la Ciencia de Datos con la atención a la estructura espacial propia de la
geografía cuantitativa. En el presente proyecto, la unidad de análisis es un lugar, la
variable de interés presenta dependencia espacial y las decisiones de escala condicionan los
resultados; por ello, el proyecto se sitúa explícitamente en el ámbito de la ciencia de
datos geoespacial y adopta sus cautelas específicas.

\subsection{Tipología del problema analítico: descriptivo, diagnóstico, predictivo y estimación transversal}
\label{subsec:tipologia-problema}

Los problemas analíticos suelen clasificarse según la pregunta que responden: el análisis
descriptivo resume lo ocurrido; el diagnóstico indaga por qué ocurrió; el predictivo estima
valores desconocidos a partir de información disponible; y el prescriptivo recomienda
acciones. Dentro del análisis predictivo conviene distinguir tres tareas que se confunden con
frecuencia. La primera es la estimación de un valor esperado condicional,
$\mu_i = E[Y_i \mid P_i, X_i]$, para unidades observadas en un mismo período, de carácter
transversal. La segunda es la proyección por escenarios, que consiste en recalcular ese valor
esperado bajo supuestos alternativos sobre las variables explicativas (por ejemplo, una nueva
distribución de la población). La tercera es el pronóstico de series temporales, que
extrapola la evolución de una variable hacia períodos futuros a partir de su dinámica pasada
\parencite{hyndman2021}.

El proyecto corresponde a la primera tarea: se estima, para cada zona, el número esperado de
consultas del período de observación dado su población y su contexto territorial, y ese
valor se contrasta con lo observado. No se pretende pronosticar cuántas consultas ocurrirán
en semanas futuras; la proyección se plantea únicamente como escenario, mediante la
actualización de la población y de los datos de consultas y la reestimación del modelo. Esta
distinción tiene una consecuencia metodológica directa: si la pregunta es cuántas consultas
cabe esperar en una zona dada su población y su contexto, la capacidad del modelo debe
evaluarse en lugares no vistos durante el entrenamiento, y no en períodos futuros. De ahí
que la validación se diseñe con particiones espaciales (véase \ref{subsec:validacion-espacial})
y no con particiones temporales.

\subsection{Datos de conteo: Poisson, exposición, sobredispersión y exceso de ceros}
\label{subsec:datos-conteo}

La variable dependiente del proyecto es un conteo, es decir, un número entero no negativo de
eventos ocurridos en una unidad espacial durante un período. El modelo de referencia para
conteos es la distribución de Poisson, cuya función de probabilidad es
\[
  P(Y = y) = \frac{e^{-\mu}\,\mu^{y}}{y!}, \qquad y = 0, 1, 2, \ldots,
\]
y cuya propiedad característica es la igualdad entre media y varianza,
$E(Y) = \operatorname{Var}(Y) = \mu$ \parencite{cameron2013}. La distribución surge
naturalmente cuando los eventos ocurren de manera independiente a una tasa constante, y
constituye el punto de partida de los modelos lineales generalizados para conteos
\parencite{mccullagh1989,agresti2015}.

Cuando las unidades difieren en tamaño, el conteo esperado es proporcional a una medida de
exposición. Si $P_i$ es la población de la zona~$i$ y $\lambda_i$ la tasa de consultas por
habitante, entonces $\mu_i = P_i \cdot \lambda_i$ y, tomando logaritmos,
$\log \mu_i = \log P_i + \log \lambda_i$. El término $\log P_i$, que entra en el modelo con
coeficiente fijo igual a uno, se denomina \textit{offset}; su inclusión garantiza que el
modelo describa tasas por habitante y que las demás variables expliquen desviaciones
respecto de la proporcionalidad con la población \parencite{cameron2013,hilbe2011}. Las
tasas suelen expresarse por cada mil habitantes, $\lambda_i \times 1000$, para facilitar su
lectura. Esta formulación es la que justifica tratar la población como exposición y no como
una variable explicativa más.

En la práctica, los conteos presentan con frecuencia una varianza superior a la media,
fenómeno denominado sobredispersión, que puede deberse a heterogeneidad no observada entre
unidades, a dependencia entre eventos o a la agregación de procesos distintos. Un
diagnóstico habitual es el estadístico de dispersión
\[
  \hat{\varphi} = \frac{1}{n - p} \sum_i \frac{(y_i - \hat{\mu}_i)^2}{\hat{\mu}_i},
\]
que en un modelo de Poisson correctamente especificado debería ser cercano a uno; valores
claramente superiores indican sobredispersión \parencite{hilbe2011}. Ignorarla no sesga
necesariamente las estimaciones puntuales, pero subestima sus errores estándar y conduce a
inferencias excesivamente confiadas \parencite{cameron2013}.

Un caso particular es el exceso de ceros, que se presenta cuando se observan más unidades
con conteo nulo de las que el modelo predice. Bajo un modelo de Poisson, el número esperado
de ceros es $\sum_i e^{-\hat{\mu}_i}$, cantidad que puede compararse con el número observado.
En el proyecto cabe esperar zonas pobladas con cero consultas, y conviene distinguir
conceptualmente entre ceros muestrales, compatibles con un valor esperado pequeño, y ceros
estructurales, que se producen porque en la zona no existe la posibilidad del evento (por
ejemplo, porque ninguna ruta mapeada la sirve y la aplicación no se usa allí). Esta
distinción motiva la consideración de modelos para sobredispersión y ceros descrita en
\ref{subsec:binomial-negativa}.

\subsection{Autocorrelación espacial: Tobler, I de Moran y LISA}
\label{subsec:autocorrelacion}

La primera ley de la geografía, formulada por \textcite{tobler1970}, establece que todo está
relacionado con todo lo demás, pero las cosas cercanas están más relacionadas que las
distantes. En términos estadísticos, esta regularidad se denomina autocorrelación espacial:
los valores de una variable en lugares próximos tienden a parecerse más (autocorrelación
positiva) o a diferir más (autocorrelación negativa) de lo que se esperaría si se
distribuyeran al azar. Su presencia viola el supuesto de independencia de las observaciones
en el que se basan muchos procedimientos estadísticos y de validación.

La medida global clásica es el índice $I$ de Moran \parencite{moran1950}, definido como
\[
  I = \frac{n}{S_0} \cdot
      \frac{\sum_i \sum_j w_{ij}\,(x_i - \bar{x})(x_j - \bar{x})}{\sum_i (x_i - \bar{x})^2},
\]
donde $n$ es el número de unidades, $x_i$ el valor en la unidad~$i$, $\bar{x}$ la media,
$w_{ij}$ el peso espacial entre las unidades $i$ y $j$, y $S_0 = \sum_i \sum_j w_{ij}$. Los
pesos se recogen en una matriz $W$ que define la vecindad. En una grilla hexagonal, la
vecindad de contigüidad resulta especialmente natural: cada zona tiene seis vecinas que
comparten con ella una arista completa (el primer anillo, obtenido en H3 con la operación
\var{grid_disk} de radio $k = 1$), de modo que $w_{ij} = 1$ si $j$ pertenece al primer
anillo de $i$ y $0$ en otro caso; habitualmente la matriz se estandariza por filas para que
cada fila sume uno. El valor esperado de $I$ bajo ausencia de autocorrelación es
$-1/(n - 1)$, y su significación se evalúa comúnmente mediante permutaciones aleatorias de
los valores entre las unidades.

Para localizar dónde se concentra la dependencia se emplean los indicadores locales de
asociación espacial (LISA) propuestos por \textcite{anselin1995}, cuya versión local del
índice de Moran es $I_i = (z_i / m_2) \sum_j w_{ij} z_j$, con $z_i = x_i - \bar{x}$ y
$m_2 = \sum_i z_i^2 / n$. La suma de los $I_i$ es proporcional al $I$ global, y cada
indicador permite clasificar las unidades en agrupamientos alto-alto y bajo-bajo, o en
atípicos alto-bajo y bajo-alto, con la debida corrección por comparaciones múltiples.
Finalmente, el correlograma espacial, que representa el valor de $I$ para vecindades de
orden creciente (anillos $k = 1, 2, 3, \ldots$), describe cómo decae la dependencia con la
distancia. En el proyecto, la autocorrelación se examina en las consultas, en las tasas y en
los residuos: el correlograma orienta la elección del tamaño de los bloques de validación
(\ref{subsec:validacion-espacial}) y el análisis de los residuos permite diagnosticar si los
modelos dejan estructura espacial sin explicar (\ref{subsec:residuos}).

\subsection{Sistemas de grillas discretas globales y H3}
\label{subsec:h3}

Un sistema de grillas discretas globales (DGGS) es una partición jerárquica de la superficie
terrestre en celdas de forma y tamaño aproximadamente regulares, a múltiples resoluciones.
\textcite{sahr2003} sistematizaron la construcción de estos sistemas geodésicos a partir de
un poliedro base, una orientación, una forma de celda, una apertura (razón de refinamiento
entre resoluciones) y una proyección entre las caras del poliedro y la esfera, y mostraron
que las grillas hexagonales basadas en el icosaedro ofrecen propiedades geométricas
favorables. H3, desarrollado originalmente por Uber y hoy mantenido como proyecto de código
abierto, es un DGGS hexagonal jerárquico construido sobre las caras de un icosaedro,
proyectadas a la esfera mediante una proyección gnomónica centrada en cada cara y orientado
de modo que sus doce vértices caen en el océano. Su resolución~0 consta de 122 celdas base
(110 hexágonos y 12 pentágonos), y cada resolución siguiente se obtiene con una apertura~7,
de modo que cada hexágono tiene siete hijos en la resolución inmediatamente más fina; en
todas las resoluciones existen exactamente 12 pentágonos, centrados en los vértices del
icosaedro \parencite{h3-overview,h3-restable}.

El hexágono presenta ventajas analíticas frente a otras teselaciones y frente a los
\textit{buffers}. A diferencia de la grilla cuadrada, en la que cada celda tiene cuatro
vecinas por arista y cuatro por vértice a distancias distintas, en la grilla hexagonal las
seis vecinas comparten una arista completa y sus centroides son equidistantes del centroide
de la celda, lo que define una vecindad uniforme y sin ambigüedad. A diferencia de los
\textit{buffers} circulares alrededor de puntos, las celdas no se solapan y cubren el
territorio de manera exhaustiva, de modo que cada evento se asigna a una sola unidad. Deben
reconocerse también sus limitaciones: la jerarquía de apertura~7 es solo aproximadamente
anidada, porque los siete hijos de un hexágono no reproducen exactamente su contorno; el
área de las celdas varía con su posición respecto de los vértices del icosaedro, con una
razón entre el hexágono de mayor y el de menor área cercana a 1,99 en la resolución~8; y los
pentágonos tienen aproximadamente la mitad del área de un hexágono de la misma resolución,
aunque, al situarse en el océano, carecen de relevancia práctica para Cochabamba
\parencite{h3-restable}.

La tabla~\ref{tab:resoluciones-h3-marco} recoge las características de las resoluciones
pertinentes para el proyecto según la tabla oficial de estadísticas de celdas de la
versión~4.x de H3. Las áreas se calculan con un modelo esférico de la Tierra basado en el
radio autálico de WGS84, y las longitudes medias de arista se calcularon exactamente para
las resoluciones 0 a 6 y se extrapolaron para las resoluciones más finas.

\begin{table}[H]
  \centering
  \caption{Resoluciones H3 relevantes para el proyecto}
  \label{tab:resoluciones-h3-marco}
  \interlineadoSimple
  \begin{tabularx}{\textwidth}{@{}>{\centering\arraybackslash}p{0.15\textwidth}
      >{\raggedleft\arraybackslash}p{0.14\textwidth}
      >{\raggedleft\arraybackslash}p{0.14\textwidth}
      >{\raggedleft\arraybackslash}p{0.16\textwidth}
      >{\raggedright\arraybackslash}X@{}}
    \toprule
    \textbf{Resolución} & \textbf{Área media del hexágono (km²)} &
    \textbf{Longitud media de arista (km)} & \textbf{Número total de celdas en el globo} &
    \textbf{Uso en el proyecto} \\
    \midrule
    6 & 36,1291 & 3,7245 & 14\,117\,882 & Macrozona (bloques de validación) \\
    7 & 5,1613 & 1,4065 & 98\,825\,162 & Análisis de sensibilidad (escala más gruesa) \\
    8 & 0,7373 & 0,5314 & 691\,776\,122 & Zona (unidad de análisis) \\
    9 & 0,1053 & 0,2008 & 4\,842\,432\,842 & Análisis de sensibilidad (escala más fina) \\
    \bottomrule
  \end{tabularx}
  \fuente{Elaboración propia a partir de \textcite{h3-restable}, versión 4.x.}
\end{table}

A partir de esta tabla se definen las dos unidades espaciales del proyecto. Se denomina zona
(o zona hexagonal) a la celda H3 de resolución~8, con un área media de aproximadamente
0,737~km² y una arista media de aproximadamente 0,531~km, lo que corresponde a un hexágono
de algo menos de un kilómetro entre lados opuestos, una escala asimilable a la de un barrio
urbano y coincidente con la resolución nativa de Kontur Population. Se denomina macrozona a
la celda H3 de resolución~6, con un área media de aproximadamente 36,13~km² y una arista
media de aproximadamente 3,72~km, esto es, un hexágono de unos 6,5~km entre lados opuestos,
cuya escala es asimilable a la de un distrito urbano. Dado que la apertura es~7, cada
macrozona agrupa aproximadamente $7^2 = 49$ zonas; la relación es aproximada porque la
jerarquía de H3 no está perfectamente anidada. Debe aclararse que, en este documento, el
término \enquote{zona} no corresponde a las zonas, distritos u organizaciones territoriales
de base (OTB) de la división administrativa municipal, sino exclusivamente a la celda H3 así
definida.

\subsection{Problema de la unidad de área modificable (MAUP)}
\label{subsec:maup}

El problema de la unidad de área modificable (MAUP, por sus siglas en inglés) designa la
dependencia de los resultados de un análisis espacial respecto de la forma en que se definen
las unidades de agregación \parencite{openshaw1984}. Se descompone en dos efectos: el efecto
de escala, por el cual los resultados cambian al variar el tamaño de las unidades (por
ejemplo, al agregar zonas en unidades mayores), y el efecto de zonificación, por el cual
cambian al variar la forma o la delimitación de unidades de tamaño similar.
\textcite{fotheringham1991} demostraron que los coeficientes de modelos de regresión
multivariante estimados con datos agregados pueden variar considerablemente, e incluso
cambiar de signo, según la escala y la zonificación empleadas, lo que los llevó a advertir
sobre la fiabilidad de los análisis multivariados basados en datos por área.

En el proyecto, la adopción de una grilla H3 regular y fijada de antemano reduce el
componente arbitrario de la zonificación, porque las celdas no se trazan en función de los
datos, pero no elimina el efecto de escala. Por ello, se recomienda, como extensión, un
análisis de sensibilidad a la resolución que repita el procedimiento con celdas de
resolución~7 (más gruesas) y~9 (más finas), y que examine si la ordenación de zonas
prioritarias y las conclusiones sustantivas se mantienen al cambiar de escala. Ese análisis
no se ejecutó en este ciclo (sección~\ref{sec:interpretacion-resultados}). Este tratamiento no resuelve el MAUP, que carece de
solución general, pero permite informar con transparencia en qué medida las conclusiones
dependen de la resolución elegida.

\subsection{Datos faltantes y discontinuidades temporales}
\label{subsec:datos-faltantes}

\textcite{rubin1976} estableció el marco de referencia para el análisis de datos faltantes,
que distingue tres mecanismos. Los datos faltan completamente al azar (MCAR) cuando la
probabilidad de ausencia no depende ni de los valores observados ni de los no observados;
faltan al azar (MAR) cuando la probabilidad de ausencia depende solo de valores observados;
y no faltan al azar (MNAR) cuando depende de los propios valores no observados. La validez
de los métodos de tratamiento depende del mecanismo: el análisis de casos completos es
insesgado bajo MCAR, la imputación múltiple es válida bajo MAR, y bajo MNAR ningún método
resulta válido sin supuestos adicionales no verificables
\parencite{little2019,vanbuuren2018}.

Los datos de consultas del proyecto presentan un período de siete semanas consecutivas sin
registros, atribuido a un cambio en el proceso de exportación. Un mecanismo de ausencia de
origen técnico, ajeno al comportamiento de los usuarios, es plausiblemente independiente de
la zona, lo que lo aproxima a MCAR en la dimensión espacial; sin embargo, no puede
descartarse que las semanas perdidas difieran sistemáticamente de las observadas en su
dimensión temporal (por ejemplo, por estacionalidad académica o festividades). La imputación
podría, en principio, rellenar esas semanas, pero sus límites son claros: la imputación
múltiple recupera la incertidumbre solo si el modelo de imputación es adecuado
\parencite{vanbuuren2018}, y aquí no existe información auxiliar específica de esas semanas
que no esté ya contenida en las semanas observadas.

Por estas razones, en el proyecto no se generan datos sintéticos para entrenar los modelos.
Hacerlo introduciría circularidad, porque los valores imputados reproducirían el patrón del
propio modelo de imputación y el modelo posterior aprendería en parte de sí mismo; crearía
falsa precisión, porque el tamaño aparente de la muestra aumentaría sin información nueva; y
podría inducir fuga de datos, si el procedimiento de imputación empleara información de
zonas destinadas a validación o prueba. En su lugar, se normaliza el conteo por el período
efectivamente observado, expresando las consultas como tasa por semana observada
($Y_i / W_{\mathrm{obs}}$, donde $W_{\mathrm{obs}}$ es el número de semanas con registros),
y se realiza un análisis de sensibilidad que compara los resultados obtenidos con los datos
anteriores y posteriores al período sin registros. Dado que el estimando es transversal, la
ausencia de semanas afecta sobre todo al nivel global de las consultas, que la normalización
corrige, y en menor medida a su distribución espacial, que el análisis de sensibilidad
permite examinar.

\subsection{Fuga de datos}
\label{subsec:fuga-datos}

\textcite{kaufman2012} definieron la fuga de datos (\textit{data leakage}) como la
introducción, en el proceso de entrenamiento, de información sobre la variable objetivo que
no estaría legítimamente disponible en el momento de usar el modelo, lo que produce
estimaciones de desempeño optimistas que no se sostienen en la aplicación real. La fuga
puede producirse por la inclusión de variables que son consecuencia del objetivo o que lo
contienen parcialmente; por la contaminación entre los conjuntos de entrenamiento y
evaluación, por ejemplo al realizar el preprocesamiento (escalado, selección de variables,
imputación) sobre todos los datos antes de particionarlos; y por la dependencia entre
observaciones de entrenamiento y evaluación que no se respeta en la partición.
\textcite{kapoor2023}, a partir de una revisión de la literatura en campos que han adoptado
el aprendizaje automático, identificaron 17 disciplinas en las que se ha detectado fuga, que
afecta colectivamente a 294 artículos y que en algunos casos condujo a conclusiones
extremadamente optimistas; los autores propusieron una taxonomía de tipos de fuga y fichas
de información del modelo para prevenirla.

En problemas espaciales existe una forma específica de fuga que resulta crítica para este
proyecto: la que se produce en predictores construidos con información de vecinos. Una
variable como la tasa de consultas de las zonas vecinas es legítima cuando se calcula con
información disponible, pero si, al evaluar el modelo en una zona de validación, la tasa
vecinal incluye los conteos de otras zonas de validación, o de la propia zona, el modelo
accede indirectamente a la variable que debe estimar. Para evitarlo, en el proyecto las
variables basadas en la vecindad se calculan excluyendo siempre la propia zona y, dentro de
cada iteración de validación, empleando únicamente los conteos de las zonas del conjunto de
entrenamiento, de modo que ninguna información del bloque evaluado intervenga en la
construcción de sus predictores.

\subsection{Sobreajuste, compromiso sesgo-varianza y validación con estructura espacial}
\label{subsec:validacion-espacial}

Se denomina sobreajuste a la situación en la que un modelo se adapta a las particularidades
aleatorias de los datos de entrenamiento en lugar de a la estructura general del fenómeno,
de modo que su desempeño en datos nuevos es inferior al observado en entrenamiento.
\textcite{hastie2009} formalizan este fenómeno mediante la descomposición del error esperado
de predicción en un punto $x_0$,
\[
  E\bigl[(Y - \hat{f}(x_0))^2\bigr] = \sigma^2 + \bigl[\mathrm{Sesgo}(\hat{f}(x_0))\bigr]^2
    + \operatorname{Var}(\hat{f}(x_0)),
\]
donde $\sigma^2$ es el error irreducible, el sesgo refleja la incapacidad del modelo para
representar la relación verdadera y la varianza su sensibilidad a la muestra de
entrenamiento. Los modelos más flexibles reducen el sesgo a costa de aumentar la varianza, y
el objetivo de la selección de modelos es encontrar el equilibrio que minimice el error en
datos no vistos, lo que exige estimar ese error de manera honesta.

La validación cruzada aleatoria, que reparte las observaciones entre particiones al azar,
presupone que las observaciones son independientes. Cuando existe autocorrelación espacial,
las observaciones de validación tienen casi siempre vecinas muy similares en el conjunto de
entrenamiento, y el error estimado refleja la capacidad del modelo para interpolar entre
vecinos más que para generalizar a lugares nuevos. \textcite{roberts2017} mostraron que, con
datos que presentan estructura temporal, espacial, jerárquica o filogenética, la validación
aleatoria produce estimaciones de error optimistas, y recomendaron la validación por
bloques, en la que las particiones se forman con grupos de observaciones contiguas, de
tamaño acorde con el alcance de la dependencia. \textcite{valavi2019} desarrollaron el
paquete blockCV para generar particiones espacial o ambientalmente separadas, y
\textcite{ploton2020} encontraron que, al validar con puntos espacialmente independientes de
los de calibración, modelos de cartografía ecológica de gran escala que parecían muy
precisos mostraban un poder predictivo muy pobre.

Este enfoque no está exento de debate. \textcite{wadoux2021} sostienen que la validación
cruzada espacial carece de fundamento teórico como método para estimar la exactitud de un
mapa y que puede resultar pesimista, porque las particiones por bloques obligan al modelo a
extrapolar a distancias mayores que las que separan los puntos a predecir de los datos
disponibles en la aplicación real; los autores defienden el muestreo probabilístico y la
inferencia basada en el diseño como vía para obtener estimaciones insesgadas.
\textcite{meyer2022} sitúan esta controversia en el contexto más amplio de la evaluación de
mapas producidos con aprendizaje automático y subrayan que la estrategia de validación debe
corresponder a la tarea de predicción: la validación aleatoria es adecuada cuando se predice
en lugares cercanos a los datos de entrenamiento, y la validación espacial, cuando se
predice en áreas alejadas de ellos.

En el proyecto, la pregunta de interés, es decir, cuántas consultas cabe esperar en una zona
dada su población y su contexto, se aplica a zonas cuya información de consultas no debería
intervenir en su propia estimación, incluidas las zonas pobladas sin consultas. Por ello, se
adopta la validación cruzada espacial por bloques con macrozonas H3 de resolución~6 como
bloques, de modo que todas las zonas de una macrozona se asignan conjuntamente a la misma
partición, y se reconoce explícitamente que esta elección puede resultar algo pesimista. El
diseño se completa con dos salvaguardas: un conjunto de prueba reservado, formado por
macrozonas completas apartadas desde el inicio, que se utiliza una sola vez al final para
evaluar el modelo elegido; y una regla de selección declarada de antemano, basada en la
devianza de Poisson media en validación cruzada y en un criterio de parsimonia, como la
regla de un error estándar descrita por \textcite{hastie2009}, que prefiere el modelo más
simple cuyo error no difiere del mínimo en más de un error estándar. Ambas salvaguardas
impiden que la elección del modelo se ajuste, consciente o inconscientemente, al conjunto de
prueba.
